\documentclass[12pt,a4paper]{letter}

\usepackage[T1]{fontenc}
\usepackage[utf8]{inputenc}
\usepackage[english]{babel}
\usepackage[margin=2.5cm]{geometry}

\begin{document}
\pagenumbering{gobble}

\begin{flushright}
Tom De Caluw\'e\\
Sonnenhofstrasse 7\\
8853 Lachen SZ\\
decaluwe.t at gmail.com\\
+41 78 448 48 48
\end{flushright}

\vspace{1em}

Hiring Team\\
\vspace{1em}

Dear Hiring Team,

I am writing to apply for the Simulacrum position on the agile Simulacrum team. The chance to work on high‑reliability instrument simulations and to improve the software stack and toolchain for medical systems is precisely the kind of engineering challenge I enjoy.

My core strengths are practical software engineering and systems thinking. I work primarily in Python for tooling, automation and backend services, and apply the same engineering practices—clear interfaces, automated tests and thorough documentation—when contributing to lower‑level code. I also have hands‑on experience in modern C++ (C++17), comfortable working on the GNU command line, and using git and collaboration platforms in team environments.

Relevant examples from my work:
- At LeeGroup I led integration work connecting e‑commerce and ERP systems; I built reliable integrations and automation in Python that reduced manual steps and improved process reliability.
- I developed an Android application integrating Zebra and NordicID RFID hardware, which involved device programming and Bluetooth‑serial communication and gave me practical experience in hardware‑near development and field troubleshooting.
- I implemented small prototypes on Arduino and Raspberry Pi for sensor and communication tasks, and have worked on multithreaded CPU and GPU workloads in other projects, which has given me experience with concurrency, performance tuning and cross‑platform debugging.

Technically I am experienced with Linux environments (daily use, shell scripting and packaging workflows), containerisation with Docker, and building reproducible development and CI workflows. I write command‑line tooling in Python and routinely use git, bash and collaborative issue/merge workflows to keep repositories and releases manageable.

I enjoy pragmatic problem solving in agile teams, collaborating closely with internal customers to turn requirements into tested, deployable software. I am also interested in Rust and modern CI practices and would welcome the opportunity to contribute to relentless improvement of the toolchain and simulation stack.

Thank you for considering my application. If you’d like, I can provide brief technical notes or a short portfolio of relevant integrations and prototypes.

Kind regards,

Tom De Caluw\'e

\end{document}
